% TOP_PROBLEM: {"problemId": "problem.bochner-riesz-conjecture", "canonicalTitle": "Bochner-Riesz conjecture", "releaseRank": 180, "primaryDomain": "analysis_pde", "primaryDomainLabel": "Analysis & PDE", "displayStatus": "open_with_solved_subcases", "releaseStatus": "open_with_solved_subcases"}
% !TeX program = lualatex
% Definition and sources only; this file is not an LLM solution attempt.
\documentclass[11pt]{article}
\usepackage[margin=25mm]{geometry}
\usepackage{fontspec}
\setmainfont{Latin Modern Roman}
\usepackage{amsmath,amssymb,mathtools}
\usepackage{unicode-math}
\setmathfont{Latin Modern Math}
\usepackage{xurl,hyperref}
\hypersetup{colorlinks=true,urlcolor=blue}
\setcounter{secnumdepth}{0}
\setlength{\parindent}{0pt}
\setlength{\parskip}{5pt}
\setlength{\emergencystretch}{3em}
\begin{document}
\section*{\texorpdfstring{Bochner-Riesz conjecture}{Bochner-Riesz conjecture}}\label{180}
\subsection{Definitions and mathematical statement}
For $R>0$, define the Euclidean Bochner--Riesz operator on Schwartz functions by
\[
 \widehat{S_R^\delta f}(\xi)=\left(1-\frac{|\xi|^2}{R^2}\right)_+^\delta\widehat f(\xi),
\]
where at $\delta=0$ the multiplier is the indicator of the ball, up to its irrelevant boundary. Uniform $L^p$ boundedness means
$\sup_{R>0}\|S_R^\delta\|_{L^p\to L^p}<\infty$.
For $d\ge2$ and $1<p<\infty$, the conjectured strong-type range is $\delta\ge0$ when $p=2$, and
$\delta>\max\{d|1/p-1/2|-1/2,0\}$ when $p\ne2$.
The problem is this full characterization in every dimension. Strong operator boundedness is distinct from weak-type endpoint estimates or almost-everywhere convergence, even when those topics appear in the source title.
\par\smallskip\textit{Formulation and concept references:} \href{https://londmathsoc.onlinelibrary.wiley.com/doi/full/10.1112/jlms.12401}{[S1]}.

\subsection{Short English statement}
For exactly which summation orders and integrability exponents do spherical Fourier cutoffs remain uniformly bounded on Euclidean Lp spaces?
\subsection{Sources}
\begin{itemize}
\item[S1]Almost everywhere convergence of Bochner-Riesz means on Heisenberg-type groups.
\url{https://londmathsoc.onlinelibrary.wiley.com/doi/full/10.1112/jlms.12401}.
\textit{Formulation source.}
\end{itemize}
\end{document}
